\Needspace{6\baselineskip}
\begin{center}\SourceLeaf{248}\pdfbookmark[2]{176.}{article-176}176.\end{center}
Ebenso wie bei beliebigen fest vorgegebenen Werten der Unbekannten jedem System von Werten der Funktionen $V$, $V'$, $V''$ usw. vor der Beobachtung eine bestimmte Wahrscheinlichkeit zukommt, so kommt umgekehrt, nachdem sich aus den Beobachtungen bestimmte Funktionswerte ergeben haben, jedem System von Werten der Unbekannten, aus dem sie hervorgegangen sein können, eine bestimmte Wahrscheinlichkeit zu: Offenbar sind diejenigen Systeme als wahrscheinlicher anzusehen, bei denen das tatsächlich eingetretene Ereignis eher zu erwarten war. Die Bestimmung dieser Wahrscheinlichkeit beruht auf folgendem Satz:

\emph{Ist unter einer Hypothese $H$ für ein bestimmtes Ereignis $E$ die Wahrscheinlichkeit $=h$, unter einer anderen Hypothese $H'$ aber, die jene ausschließt und an sich ebenso wahrscheinlich ist, die Wahrscheinlichkeit desselben Ereignisses $=h'$: Dann behaupte ich, dass nach dem tatsächlichen Eintreten des Ereignisses $E$ die Wahrscheinlichkeit, dass $H$ die wahre Hypothese war, sich zur Wahrscheinlichkeit, dass $H'$ die wahre Hypothese war, verhält wie $h$ zu $h'$}\HistoricalNoteEight{Handschriftliche Bemerkung:] Hätten die Hypothesen $H$, $H'$ an sich (d.\,h. vor dem Eintreten von $E$ oder bevor dieses Eintreten bekannt wurde), ungleiche Wahrscheinlichkeiten $\mu$, $\mu'$ gehabt, so muss man ihnen, nach dem Eintreten von $E$, Wahrscheinlichkeiten zuweisen, die den Produkten $\mu h$, $\mu' h'$ proportional sind.}.

Zum Beweis nehmen wir an, dass wir, indem wir alle Umstände unterscheiden, von denen es abhängt, ob $H$ oder $H'$ oder eine andere Hypothese zutrifft und ob das Ereignis $E$ oder ein anderes eintreten muss, ein System verschiedener Fälle erhalten. Jeder dieser Fälle ist an sich (d.\,h. solange ungewiss ist, ob das Ereignis $E$ oder ein anderes eintreten wird) als gleich wahrscheinlich anzusehen. Diese Fälle sollen sich wie folgt verteilen,
\begin{center}
\small
\begin{tabular}{c|c|c}
\shortstack{sodass darunter\\so viele Fälle sind} & \shortstack{in denen die Hypothese\\zutreffen muss} & \shortstack{unter solchen Bedingungen,\\dass das Ereignis sein muss} \\
\hline
$m$ & $H$ & $E$ \\
$n$ & $H$ & von $E$ verschieden \\
$m'$ & $H'$ & $E$ \\
$n'$ & $H'$ & von $E$ verschieden \\
$m''$ & von $H$ und $H'$ verschieden & $E$ \\
$n''$ & von $H$ und $H'$ verschieden & von $E$ verschieden
\end{tabular}
\end{center}
Dann gilt $h=\dfrac{m}{m+n}\mathpunct{\raisebox{.45ex}{,}}$ $h'=\dfrac{m'}{m'+n'}\mathpunct{\raisebox{.45ex}{;}}$ ferner hatte, bevor das Ereignis bekannt wurde, die Wahrscheinlichkeit der \SourceLeaf{249}Hypothese $H$ den Wert $=\dfrac{m+n}{m+n+m'+n'+m''+n''}\mathpunct{\raisebox{.45ex}{,}}$ nachdem das Ereignis bekannt ist und die Fälle $n$, $n'$, $n''$ aus der Menge der möglichen Fälle ausscheiden, ergibt sich für die Wahrscheinlichkeit derselben Hypothese der Wert $=\dfrac{m}{m+m'+m''}\mathpunct{\raisebox{.45ex}{;}}$ ebenso ergibt sich die Wahrscheinlichkeit der Hypothese $H'$ vor und nach dem Ereignis jeweils aus $\dfrac{m'+n'}{m+n+m'+n'+m''+n''}$ und $\dfrac{m'}{m+m'+m''}\mathpunct{:}$ Da wir also voraussetzen, dass die Hypothesen $H$ und $H'$ gleich wahrscheinlich sind, bevor das Ereignis bekannt ist, gilt $m+n=m'+n'$, woraus die Behauptung unmittelbar folgt.

Soweit wir nun annehmen, dass außer den Beobachtungen $V=M$, $V'=M'$, $V''=M''$ usw. keine weiteren Daten zur Bestimmung der Unbekannten vorliegen, und somit alle Systeme von Werten dieser Unbekannten vor diesen Beobachtungen gleich wahrscheinlich waren, ist die Wahrscheinlichkeit jedes bestimmten Systems nach diesen Beobachtungen offenbar proportional zu $\Omega$ selbst. Das ist so zu verstehen, dass die Wahrscheinlichkeit, mit der die Werte der Unbekannten jeweils zwischen den unendlich nahen Grenzen $p$ und $p+\mathrm{d}p$, $q$ und $q+\mathrm{d}q$, $r$ und $r+\mathrm{d}r$, $s$ und $s+\mathrm{d}s$ usw. liegen, durch $\lambda\Omega\,\mathrm{d}p\,\mathrm{d}q\,\mathrm{d}r\,\mathrm{d}s$ usw. ausgedrückt wird, wobei $\lambda$ eine von $p$, $q$, $r$, $s$ usw. unabhängige Konstante ist. Und zwar ist $\dfrac{1}{\lambda}$ offenbar der Wert des Integrals $\nu^{ti}$ Ordnung: $\int^{\nu}\Omega\,\mathrm{d}p\,\mathrm{d}q\,\mathrm{d}r\,\mathrm{d}s\, .\, .\, .\, .$, wobei jede der Variablen $p$, $q$, $r$, $s$ usw. von $-\infty$ bis $+\infty$ läuft.
